\documentclass[10pt]{article}
\usepackage[utf8]{inputenc}
\usepackage[T1]{fontenc}
\usepackage{amsmath}
\usepackage{amsfonts}
\usepackage{amssymb}
\usepackage[version=4]{mhchem}
\usepackage{stmaryrd}

\begin{document}
2.18 Definition Let $X$ be a metric space. All points and sets mentioned below are understood to be elements and subsets of $X$.\\
(a) A neighborhood of $p$ is a set $N_{r}(p)$ consisting of all $q$ such that $d(p, q)<r$, for some $r>0$. The number $r$ is called the radius of $N_{r}(p)$.\\
(b) A point $p$ is a limit point of the set $E$ if every neighborhood of $p$ contains a point $q \neq p$ such that $q \in E$.\\
(c) If $p \in E$ and $p$ is not a limit point of $E$, then $p$ is called an isolated point of $E$.\\
(d) $E$ is closed if every limit point of $E$ is a point of $E$.\\
(e) A point $p$ is an interior point of $E$ if there is a neighborhood $N$ of $p$ such that $N \subset E$.\\
(f) $E$ is open if every point of $E$ is an interior point of $E$.\\
(g) The complement of $E$ (denoted by $E^{c}$ ) is the set of all points $p \in X$ such that $p \notin E$.\\
(h) $E$ is perfect if $E$ is closed and if every point of $E$ is a limit point of $E$.\\
(i) $E$ is bounded if there is a real number $M$ and a point $q \in X$ such that $d(p, q)<M$ for all $p \in E$.\\
(j) $E$ is dense in $X$ if every point of $X$ is a limit point of $E$, or a point of $E$ (or both).\\
Let us note that in $R^{1}$ neighborhoods are segments, whereas in $R^{2}$ neighborhoods are interiors of circles.\\
2.19 Theorem Every neighborhood is an open set.

Proof Consider a neighborhood $E=N_{r}(p)$, and let $q$ be any point of $E$. Then there is a positive real number $h$ such that

$$
d(p, q)=r-h
$$

For all points $s$ such that $d(q, s)<h$, we have then

$$
d(p, s) \leq d(p, q)+d(q, s)<r-h+h=r
$$

so that $s \in E$. Thus $q$ is an interior point of $E$.\\
2.20 Theorem If $p$ is a limit point of a set $E$, then every neighborhood of $p$ contains infinitely many points of $E$.

Proof Suppose there is a neighborhood $N$ of $p$ which contains only a finite number of points of $E$. Let $q_{1}, \ldots, q_{n}$ be those points of $N \cap E$, which are distinct from $p$, and put

$$
r=\min _{1 \leq m \leq n} d\left(p, q_{m}\right)
$$


\end{document}